%% kapa: exam=1 round=31 year=2020 period=1 form=A subject=economics subject_ko=경제학원론 seq=11 number=51
%% source: Q-Net 감정평가사 기출문제 (cst00309) artlSeq=5207219 file=extracted/1008/5207219/A형/A형/2020년도 제31회 감정평가사 1차 시험_1교시_A형.pdf
%% conversion: kordoc 4.24.51 -> tools/gen1.py (auto); review=proofread-sonnet; flags=none
\begin{question}{51}
그림과 같이 완전경쟁시장이 독점시장으로 전환되었다. 소비자로부터 독점기업에게 이전되는 소비자잉여는? (단, MR은 한계수입, MC는 한계비용, D는 시장수요곡선으로 불변이다. 독점기업은 이윤극대화를 추구한다.)
\begin{center}
\begin{tikzpicture}[x=0.1mm,y=-0.1mm,font=\scriptsize,line width=0.7pt]
\node at (208,25) {\textbf{완전경쟁시장}};
\node at (690,25) {\textbf{독점시장}};
\node[anchor=west] at (12,52) {\textbf{가격, 비용}};
\node[anchor=west] at (462,52) {\textbf{가격, 비용}};
% left panel
\draw[-{Stealth[length=1.8mm]}] (65,385) -- (65,85);
\draw[-{Stealth[length=1.8mm]}] (65,385) -- (435,385);
\draw (65,152) -- (353,358);
\draw (65,342) -- (400,342);
\draw[dashed] (333,342) -- (333,385);
\node at (195,215) {\textbf{D}};
\node at (362,322) {\textbf{MC}};
\node at (40,152) {\textbf{70}};
\node at (40,342) {\textbf{10}};
\node at (47,398) {\textbf{0}};
\node at (333,414) {\textbf{100}};
\node at (418,410) {\textbf{수량}};
% arrow
\draw[-{Stealth[length=3mm]},line width=2.2pt] (413,203) -- (442,203);
% right panel
\draw[-{Stealth[length=1.8mm]}] (513,385) -- (513,85);
\draw[-{Stealth[length=1.8mm]}] (513,385) -- (885,385);
\draw (513,152) -- (803,358);
\draw (513,152) -- (675,362);
\draw (513,342) -- (850,342);
\draw[dashed] (513,255) -- (660,255) -- (660,385);
\draw[dashed] (782,342) -- (782,385);
\node at (645,218) {\textbf{D}};
\node at (610,322) {\textbf{MR}};
\node at (812,322) {\textbf{MC}};
\node at (488,152) {\textbf{70}};
\node at (488,255) {\textbf{40}};
\node at (488,342) {\textbf{10}};
\node at (496,398) {\textbf{0}};
\node at (660,414) {\textbf{50}};
\node at (782,414) {\textbf{100}};
\node at (868,410) {\textbf{수량}};
\end{tikzpicture}
\end{center}
\begin{choices}
  \item 0
  \item 750
  \item 1,500
  \item 2,250
  \item 3,000
\end{choices}
\end{question}
\answer{3}
